\chapter{Introduction}
\label{ch:introduction}

The name \emph{Sentinel} is derived from the Old Italian word
\emph{sentinella}, referring to a person or thing tasked with watching
over, protecting, and alerting others of potential danger. This is
precisely what the device aims to do for a climber: it continuously
watches over the climbing rope, manages it to protect the climber during
the ascent, and reacts immediately -- by braking -- the moment a danger
(a fall) is detected.

\section{Background}
\label{sec:background}

Rock and sport climbing rely fundamentally on the belay system: the
mechanism and human partner responsible for managing rope slack and
arresting a climber's fall. Traditional belaying requires a trained,
attentive human belayer holding the rope at all times, which introduces
risks related to human error, fatigue, distraction, and inconsistent
reaction times. Semi-automated mechanical devices assist with braking by
using a cam or pinching mechanism that clamps the rope against a grooved
wheel, but they still require a human operator to feed rope in and out and
to manage slack during the climb.

As an initial mechanical design reference, this project looks specifically
at the Petzl GriGri, shown in \cref{fig:petzl_grigri}. Unlike a plain
friction hitch or a static cam, the GriGri already relies on a mechanism
with moving parts -- a camming/pinching lever that clamps the rope
against a grooved wheel to generate braking friction -- which makes it a
plausible starting point for turning into a mechatronic (i.e.\
motor-actuated and sensor-driven) system, rather than having to design a
rope-gripping mechanism from scratch. A free 3D CAD model of this
mechanism, available on the GrabCAD model repository~\cite{grabcad_petzl_grigri},
was used to study its mechanical layout and is used here purely as a
starting point for the mechanical design of Sentinel's own rope-gripping
mechanism, to be extended with active sensing and actuation as described
in this thesis.

\begin{figure}[H]
    \centering
    \includegraphics[width=0.55\textwidth]{figures/petzl_grigri_reference.png}
    \caption{Petzl GriGri mechanical assembly. CAD model obtained from the
    GrabCAD model repository~\cite{grabcad_petzl_grigri}, used as a
    starting point for the design of Sentinel's rope-gripping mechanism.}
    \label{fig:petzl_grigri}
\end{figure}

\section{Project Area}
\label{sec:project-area}

This thesis sits at the intersection of mechatronics, embedded real-time
systems, and sports safety equipment engineering. The project area covers:
the mechanical design of a rope-handling and braking mechanism; the
integration of motion, position, and force sensing to estimate a climber's
state in real time; the design of an active motor-driven mechanism for
slack management and emergency braking; and the analysis of relevant
safety regulations and standards for climbing protection equipment.

\section{Narrowing}
\label{sec:narrowing}

Given the breadth of the project area, this work narrows its focus to a
single-rope, indoor, controlled-environment prototype. The system is
narrowed to:

\begin{itemize}
    \item A single, standard dynamic climbing rope diameter and type,
        rather than supporting the full range of ropes on the market.
    \item Bench-top and controlled indoor testing, rather than outdoor or
        production climbing-gym deployment.
    \item A functional proof-of-concept prototype, rather than a
        manufacturable, certified consumer product.
\end{itemize}

\section{Principal Objectives}
\label{sec:principal-objectives}

\begin{enumerate}
    \item Create a functional prototype of an automated self-belaying
        device.
    \item Perform comprehensive testing and evaluation of the prototype.
    \item Analyze regulations, standards, and safety requirements
        applicable to such climbing devices.
\end{enumerate}

\section{Secondary Objectives}
\label{sec:secondary-objectives}

\begin{enumerate}
    \item Survey existing commercial and academic self-belay and
        assisted-braking solutions.
    \item Formulate comprehensive functional and technical requirement
        specifications for the system.
\end{enumerate}

\section{Research Questions}
\label{sec:research-questions}

\begin{enumerate}
    \item[\textbf{RQ1}] To what extent can an embedded sensing system
        reliably distinguish between a climber pulling in rope during
        normal ascent and the onset of an actual fall?
    \item[\textbf{RQ2}] How can an automated belay mechanism safely
        transition out of a braked state once the climber has resumed a
        valid climbing position, without requiring manual intervention?
    \item[\textbf{RQ3}] Which existing regulations and standards
        applicable to climbing protection equipment are relevant to an
        automated belay device, and what safety requirements do they
        impose on its design?
    \item[\textbf{RQ4}] What mechanical and control-system safeguards are
        necessary to guarantee a safe fallback behaviour (e.g.\ automatic
        rope locking) in the event of a power or electronic failure?
\end{enumerate}

\section{Limitations}
\label{sec:limitations}

The following are explicitly outside the scope of this thesis, unless
revisited as future work:

\begin{itemize}
    \item Formal certification of the device against climbing safety
        standards (e.g.\ UIAA, CE/EN).
    \item Outdoor field testing with real climbers at height.
    \item Support for multiple rope diameters, types, or conditions
        (e.g.\ wet or icy rope).
    \item Weatherproofing and long-term durability testing of the
        prototype enclosure.
    \item Mass-manufacturability and industrial design for a commercial
        product.
\end{itemize}

\section{Contributions}
\label{sec:contributions}

This thesis is expected to contribute:

\begin{enumerate}
    \item A working proof-of-concept prototype demonstrating that
        automated, sensor-driven belaying with real-time slack management
        and fast emergency braking is technically feasible using
        accessible embedded hardware.
    \item An analysis of existing commercial and academic self-belay
        solutions, situating this project's approach relative to the
        state of the art.
    \item A documented functional and technical requirement
        specification for automated belay devices.
    \item An analysis of applicable regulations, standards, and safety
        requirements relevant to automated belay device design.
    \item Empirical performance data (e.g.\ reaction time, braking
        distance, false-trigger rate) characterising the prototype's
        safety-relevant behaviour under controlled test conditions.
\end{enumerate}

\section{Report Structure}
\label{sec:report-structure}

This report is organised to build from problem definition to
implementation and evaluation, following the structure of a typical
software/mechatronic development-project thesis. \Cref{ch:introduction}
(this chapter) introduces the motivation, objectives, research questions,
and scope. \Cref{ch:functional-requirements} defines the functional and
non-functional requirements that the system must satisfy, at both a high
and low level. \Cref{ch:technical-design} describes the overall
architecture of the chosen solution and why it was chosen.
\Cref{ch:development-process} describes the development process
followed, including tools for project management and documentation.
\Cref{ch:implementation} covers the technical details of the
implementation (tools, languages, libraries, and representative code or
design examples). \Cref{ch:deployment} describes how the system is set
up and deployed on the target hardware. \Cref{ch:testing-feedback}
describes how the system was tested and what was learned from testing.
\Cref{ch:discussion} discusses the process and product choices made
throughout the project, and \Cref{ch:conclusion} summarises what was
achieved and outlines future work. Every chapter beyond this Introduction
and the Requirements chapter is currently an empty placeholder, to be
filled in as the project progresses; each structural change to this
document is logged in the project's development log
(\texttt{docs/DEV\_LOG.md}).

\section{Report Tools}
\label{sec:report-tools}

This section documents the tools used in preparing this report and the
underlying project work.

\subsection{AI Tools}

An AI-assisted coding and documentation environment (Cursor) was used
throughout the project to accelerate firmware scaffolding, documentation
drafting, and LaTeX authoring. All AI-assisted contributions were
reviewed and directed by the author; ambiguous design decisions were
resolved by the author rather than left to automated assumptions, per the
project's development workflow.

\subsection{CAD Tools}

Mechanical design and reference modelling work uses SolidWorks, alongside
publicly available reference CAD models (see \cref{sec:background}) used
for early-stage design inspiration.

\subsection{LaTeX and Documentation Tools}

This report is written in \LaTeX{}, using the community \texttt{ntnuthesis}
document class for the long-form thesis, compiled locally with a MiKTeX
distribution (\texttt{pdflatex}/\texttt{biber}). General project
documentation (hardware specifications, pinouts, and development history)
is maintained in Markdown under \texttt{docs/}.

\subsection{Software Development Environments}

Embedded firmware is developed in C++ using the PlatformIO build system
and toolchain, targeting an embedded microcontroller platform. Version
control is managed with Git.
